\documentclass{beamer}
\usepackage{udec-slides}
\ayud{1}{Circuitos Serie/Paralelo}

\usefonttheme{serif}

\usepackage{newtxtext,newtxmath}

\begin{document}
\primeraDiapo

% -------------------------------------------------------------
% SLIDE 1: EJERCICIO 1
% -------------------------------------------------------------
\begin{frame}{Ejercicio 1}
    \begin{figure}[htbp]
        \centering
        \begin{circuitikz}[american, scale=0.9, transform shape]
            % Fuente de voltaje
            \draw (0,0) to[battery1, l=\qty{60}{\V}, invert] (0,3);

            % Rama superior: Amperímetro y resistencia de 12 ohms
            \draw (0,3) to[ammeter, l=\qty{3.048}{\A}] (2.5,3)
            to[R, l=\qty{12}{\ohm}, -*] (5,3);

            % Rama inferior: Resistencia de 4 ohms
            \draw (0,0) to[R, l=\qty{4}{\ohm}, -*] (5,0);

            % Rama vertical central: Resistencia de 5 ohms
            \draw (5,3) to[R, l=\qty{5}{\ohm}] (5,0);

            % Malla derecha
            \draw (5,3) to[R, l=\qty{3}{\ohm}] (8,3)
            to[R, l=\qty{5}{\ohm}] (8,0)
            to[R, l=\qty{6}{\ohm}] (5,0);
        \end{circuitikz}
    \end{figure}
\end{frame}

% -------------------------------------------------------------
% SLIDE 2: EJERCICIO 2
% -------------------------------------------------------------
\begin{frame}{Ejercicio 2}
    \begin{figure}[htbp]
        \centering
        \begin{circuitikz}[american, scale=0.85, transform shape]
            % Definición de coordenadas principales
            \coordinate (top_node) at (4.5,3);
            \coordinate (bot_left) at (3,0);
            \coordinate (bot_right) at (6,0);

            % Fuente de voltaje
            \draw (0,0) to[battery1, l=\qty{25}{\V}, invert] (0,3);

            % Conexión superior con amperímetro
            \draw (0,3) to[ammeter, l=\qty{2.695}{\A}] (top_node);

            % Conexión inferior con resistencia de 6 ohms
            \draw (0,0) to[R, l=\qty{6}{\ohm}] (bot_left);

            % Configuración Delta central (triángulo)
            \draw (bot_left) to[R, l=\qty{6}{\ohm}] (top_node);
            \draw (top_node) to[R, l=\qty{5}{\ohm}] (bot_right);
            \draw (bot_left) to[R, l=\qty{4}{\ohm}] (bot_right);

            % Malla derecha
            \draw (top_node) to[R, l=\qty{3}{\ohm}] (8.5,3)
            to[R, l=\qty{4}{\ohm}] (8.5,0)
            to[R, l=\qty{2}{\ohm}] (bot_right);

            % Puntos de unión (nodos negros)
            \node[circ] at (top_node) {};
            \node[circ] at (bot_left) {};
            \node[circ] at (bot_right) {};
        \end{circuitikz}
    \end{figure}
\end{frame}

% -------------------------------------------------------------
% SLIDE 3: EJERCICIO 3
% -------------------------------------------------------------
\begin{frame}{Ejercicio 3}
    \begin{figure}[htbp]
        \centering
        \begin{circuitikz}[american, scale=0.85, transform shape]
            % Definición de coordenadas principales
            \coordinate (mid_left) at (3,2);
            \coordinate (bot_left) at (3,0);
            \coordinate (mid_right) at (8.5,2);

            % Fuente de voltaje
            \draw (0,0) to[battery1, l=\qty{75}{\V}, invert] (0,2);

            % Amperímetro hacia el nodo central-izquierdo
            \draw (0,2) to[ammeter, l=\qty{7.471}{\A}] (mid_left);

            % Resistencia inferior izquierda
            \draw (0,0) to[R, l=\qty{6}{\ohm}] (bot_left);

            % Resistencia vertical intermedia
            \draw (mid_left) to[R, l=\qty{6}{\ohm}] (bot_left);

            % Rama superior
            \draw (mid_left) to[R, l=\qty{2}{\ohm}] (3,4)
            to[R, l=\qty{4}{\ohm}] (8.5,4)
            to[R, l=\qty{5}{\ohm}] (mid_right);

            % Rama central horizontal
            \draw (mid_left) to[R, l=\qty{3}{\ohm}] (mid_right);

            % Rama inferior derecha
            \draw (mid_right) to[R, l=\qty{4}{\ohm}] (8.5,0)
            to[R, l=\qty{2}{\ohm}] (6,0)
            to[R, l=\qty{4}{\ohm}] (bot_left);

            % Puntos de unión (nodos negros)
            \node[circ] at (mid_left) {};
            \node[circ] at (bot_left) {};
            \node[circ] at (mid_right) {};
        \end{circuitikz}
    \end{figure}
\end{frame}

\end{document}
